% FILE: artifacts/_common/style/cards.tex
% ==================================================================================================
% THE CARD GRAMMAR of the house style: the tcolorbox card styles and their glow, the two builders that
% keep a tcolorbox out of a TikZ node, and the pointer marks (literature citation, supplement badge,
% QR pointer). A document inputs it after the colours and the fonts, and loads tcolorbox before it.
% \seebottom also reads \icQR from the document's icon registry and \iconlg from its type scale.

% ===== Card (panel) styles ========================================================================
% THE CARD GRAMMAR, and both halves of it are load-bearing:
%   NO BANNERS. A card's identity rides on a solid accent DISC + an accent-coloured TITLE on the plain
%   surface. A full-bleed accent band is the loudest device a card has, and it would be spent on the
%   card's NAME while the payoff below stays plain text. Do not reintroduce `colbacktitle'.
%   RADIUS IS A RATIO, NOT A NUMBER: ~4-5% of the card's own width (24 mm at 65 cm; 14 mm at 30 cm for
%   a nested card). Below ~3% a poster-scale arc reads as a chamfer. SCALE IT WITH THE CARD.
% CARD PADDING, HELD AS PLAIN cm NUMBERS. Two consumers need the same value in two forms: a tcolorbox
% key wants a dimension (`left=\gPadSideCm cm') and the region map wants a number it can do arithmetic
% on (the medallion's clearance against a question card's icon tile is derived from \gPadTopCm). One
% number serving both is what keeps the check measuring the point the layout actually places.
% TWO VERTICAL CONVENTIONS, deliberately distinct and each named: the plain `cardbase' pads 1.1 top and
% bottom, while a COLUMN card pads 1.3 / 1.4 because it carries a heading and a figure rather than a
% figure alone. Every measured box height in the poster is derived from one of these, so they are not
% interchangeable -- collapsing them means re-deriving \hBand and every row top.
\def\gPadSideCm{1.2}    % left / right, every card
\def\gPadBaseCm{1.1}    % top / bottom of a plain `cardbase' card
\def\gPadTopCm{1.3}     % top of a column card
\def\gPadBotCm{1.4}     % bottom of a column card
\tcbset{
  % shadow-free base: surface, rounded corners, padding.
  cardbase/.style={
    enhanced, arc=24mm, outer arc=24mm,
    colback=cardground, colframe=cardring, boxrule=0.8pt,
    boxsep=0pt, left=\gPadSideCm cm, right=\gPadSideCm cm,
    top=\gPadBaseCm cm, bottom=\gPadBaseCm cm,
    fontupper=\color{textprimary}\bodysize,
  },
  % ===== ONE SHADOW FOR EVERY CARD ON THE POSTER ==================================================
  % Every card glows in its identity colour through the SAME blur, cast and spread (`posterglow' below
  % -- the single copy of those numbers). Only the tint changes. TWO TRAPS, BOTH SILENT:
  %   (1) THE PRIMITIVE MATTERS AS MUCH AS THE NUMBERS. tcolorbox's `fuzzy shadow' is a stack of offset
  %       rounded rectangles; TikZ's `blur shadow' (shadows.blur) is a true gaussian blur, and giving a
  %       fuzzy shadow these numbers renders a stepped halo. Hence `blur shadow' on the FRAME PATH --
  %       the outermost path, drawn beneath the interior, so the shadow lands behind the whole card.
  %   (2) THE TINT MUST GO THROUGH `every shadow'. A `fill=' inside `blur shadow={...}' LEAKS onto the
  %       path it is attached to, and the card comes out flooded with the tint.
  %   #1 = ring colour (the frame path must keep filling itself with it)   #2 = tint of the glow
  cardglow/.style 2 args={
    frame style={fill=#1, posterglow, every shadow/.style={fill=#2, draw=none}}},
  % Section cards glow in their IDENTITY colour -- the heading, number disc and icon carry it -- and not
  % in their ring, which is a near-white hairline whose glow would be invisible. For a neutral cast:
  % {bordersubtle}{cardshadow}. The tint is `cardcast', so it follows the identity colour.
  card/.style={cardbase, cardglow={cardring}{cardcast}},
}

% A tcolorbox typeset directly inside a TikZ \node triggers an infinite NFSS font-selection loop under
% fontspec/LuaLaTeX. The fix is two-step: build the card into a savebox at document top level (OUTSIDE
% any tikzpicture), then place that box with \node{\usebox} inside the overlay picture.
% Build card CONTENT into <savebox> (call at document top level):
%   \buildcard{<\savebox>}{<width cm>}{<height cm>}{<tcb style>}{<content>}
\newcommand{\buildcard}[5]{%
  \sbox#1{\begin{tcolorbox}[#4, width=#2 cm, height=#3 cm]#5\end{tcolorbox}}%
}
% Build arbitrary (card-free) content into <savebox> at document top level:
%   \buildbox{<\savebox>}{<content>}
\newcommand{\buildbox}[2]{\sbox#1{#2}}

% ===== Inline citation + supporting-detail reference tag ==========================================
% \litcite and \seebottom share the CHROME accent because they are the same kind of mark: a pointer
% AWAY from the poster (to the literature, to the supplement). Deliberately not two colours.
\newcommand{\litcite}[1]{{\color{pointermark}\notesize (#1)}}
% Small numbered badge matching the numbered points in supplement.tex. It is placed with \raisebox (a
% LaTeX box lift) around a picture whose baseline is pinned to the node CENTRE via a node anchor --
% neither depends on the surrounding picture's axes, so the badge aligns identically in ordinary text
% and inside a flipped-axis node (e.g. x=1cm y=-1cm).
\newcommand{\refnum}[1]{\raisebox{3.2mm}{\tikz[baseline=(rn.center)]{\node(rn)[fill=pointermark,
  text=white, rounded corners=3mm, minimum size=12mm, inner sep=0pt,
  font=\headingfont\notesize]{#1};}}}
% Pointer to numbered item(s) in the standalone supplement: the QR glyph followed by the item numbers,
% signalling "scan the header QR to reach supplemental item N". A BARE GLYPH SAYS ONLY `there is more
% somewhere'; the numbers are what make it a pointer, and they are the numbers the supplement's own
% \suphead badges carry -- keep the two in step when a supplement point is added or renumbered.
% Takes a COMMA LIST, so one macro serves one, two or three targets: \seebottom{4,5,6}.
% The glyph is set at \iconlg to match the heading it sits in; the badges bring their own size.
\newcommand{\seebottom}[1]{{\color{pointermark}%
  \mbox{{\iconlg\icQR}\;\foreach \sbn [count=\sbi] in {#1}{\ifnum\sbi>1\,\fi\refnum{\sbn}}}}}

% ----- THE CARD GLOW'S THREE NUMBERS ---------------------------------------------------------------
% \providecommand, not \def, so that a document loading the colours first may pre-empt them.
% >>> WEIGHT MOVES WITH THE IDENTITY COLOUR <<< The glow is tinted with `cardcast', i.e. with the
% identity colour, and it is the accent's largest surface on the page. Under the deep petrol accent a
% strong cast reads as a bruise rather than a halo, so the glow is light and short: it reads as CONTACT
% rather than elevation. Retune these three with `chromecore'. Only `posterglow' reads them --
% `softshadow' is neutral depth INSIDE a figure.
\providecommand\gGlowDrop{-1.8mm}    % how far the glow falls below the card
\providecommand\gGlowBlur{4.5mm}     % its blur radius
\providecommand\gGlowOpacity{15}     % its strength, 0-100
\tikzset{
  % ---- shadows -------------------------------------------------------------------------------------
  % Deep glow: the poster's card / banner shadow, used by `cardglow' above and by the mechanism tree's
  % stage banners. The TINT is NOT here: it rides on `every shadow' at the call site (trap note above).
  posterglow/.style={blur shadow={shadow xshift=0pt, shadow yshift=\gGlowDrop,
    shadow blur radius=\gGlowBlur, shadow blur steps=12, shadow opacity=\gGlowOpacity}},
}
